\documentclass[../../../main.tex]{subfiles}

\begin{document}

\chapter{Vector-Valued Functions and Parametrization}

In our previous note, we reviewed about vectors and their geometry.
For calculus in higher dimension, we also study vector-valued
functions in addition to real-valued ones that we discussed for
calculus in 2-space. In this note, we will further extend our
discussion on vectors with vector-valued functions, parametric
curves, and calculus in higher dimension.

\section{Vector-Valued Functions and Calculus}

As always, let's get started with the definition.

\begin{definition}{}{}
A function with real inputs is a \textbf{vector-valued function}
if its outputs are vectors. The \textbf{implicit domain} is the
largest set of reals for which the function is well-defined
when the domain of a vector-valued function is not explicitly
stated.
\end{definition}

Below are some examples of vector-valued functions.
\begin{align*}
    \mathbf{F}(t) &= \langle t + 2, 3t^2, 4t^3 \rangle,
        \quad -10 \leq t \leq 10 \\
    \mathbf{r}(t) &= t \mathbf{i} + 5t \mathbf{j} + \cos t \mathbf{k},
        \quad -2\pi \leq t \leq 10\pi
\end{align*}
Note that the range need not be vectors in 3-space, but it could
be any vectors. From the example above, we could define another term.

\begin{definition}{}{}
Consider a vector-valued function $\mathbf{r} = f_1 (t) \mathbf{i}
+ f_2 (t) \mathbf{j} + f_3 (t) \mathbf{k}$. The functions $f_1, f_2,
f_3$ of $t$ are known as the \textbf{component functions} of
$\mathbf{r}$.
\end{definition}

Graphing such functions is pretty intuitive. For each input $t$,
we would draw the resulting vector from the origin, which is called
the \textbf{radius vector} or \textbf{position vector}. Below is an
example of graphing $\mathbf{r} = \cos t \mathbf{i} +
\sin t \mathbf{j} + t \mathbf{k}$ for $t \in [-\pi, 2\pi]$.

\begin{figure}[H]
\centering
\begin{tikzpicture}[scale=1.4]
\begin{axis}[
    view={120}{35},
    axis lines=center,
    grid=none,ticks=none,
    samples=300,
    samples y=0,
    xlabel={$x$},ylabel={$y$},zlabel={$z$},
    xmin=-2, xmax=2,
    ymin=-2, ymax=2,
    zmin=-3.5, zmax=6.5
]
    \pgfmathsetmacro{\coorda}{cos(deg(pi))}
    \pgfmathsetmacro{\coordb}{sin(deg(pi))}
    \pgfmathsetmacro{\coordc}{pi}
    \pgfmathsetmacro{\coordd}{cos(deg(pi/2-0.4))}
    \pgfmathsetmacro{\coorde}{sin(deg(pi/2-0.4))+0.05}
    \pgfmathsetmacro{\coordf}{pi/2-0.4}
    \pgfmathsetmacro{\coordg}{cos(deg(pi/2+0.4))}
    \pgfmathsetmacro{\coordh}{sin(deg(pi/2+0.4))+0.05}
    \pgfmathsetmacro{\coordi}{pi/2+0.4}

    \addplot3[
        very thick,
        Red,
        domain=-pi:2*pi,
        variable=\t,
    ] ({cos(deg(\t))}, {sin(deg(\t))}, {\t});

    \addplot3[very thick,Blue,-{Stealth[length=2mm]}]
    coordinates {
        (0,0,0)
        (\coorda,\coordb,\coordc)
    };

    \addplot3[line width=4pt,Green,-{Stealth[length=2mm,width=4mm]}]
    coordinates {
        (\coordd,\coorde,\coordf)
        (\coordg,\coordh,\coordi)
    };
\end{axis}
\end{tikzpicture}
\end{figure}

As you can see from the figure above, the blue arrow from the origin
is a radius vector of the helix. The arrowhead on the curve is known
as the \textbf{orientation} of the graph, which represents the
direction in which the curve traverse as the input increase.

\begin{important}
One more important thing to note is that the parametrization of a
curve is \textit{not} unique. For instance, the curve above
$\mathbf{r} = \cos t \mathbf{i} + \sin t \mathbf{j} + t \mathbf{k}$
for $t \in [-\pi, 2\pi]$ can also be represented as
$\mathbf{s} = \cos 2t \mathbf{i} + \sin 2t \mathbf{j} + 2t \mathbf{k}$
for $t \in [-\frac{\pi}{2}, \pi]$.
\end{important}

Extending our example, we can see that the curve also lies on an
circular curve since for $x(t) = \cos t$ and $y(t) = \sin t$,
$x^2 + y^2 = 1$.

\subsection{Calculus for Vector-Valued Functions}

Now that we discussed what vector-valued functions are and how to
graph them, it is natural to see how calculus work for such functions.
Just like graphing, it is quite natural to see how calculus works
given the properties of limits, derivatives, and integrals.

\begin{important}
Given a vector-valued function $\mathbf{r}(t) = f_1(t) \mathbf{i}
+ f_2(t) \mathbf{j} + f_3(t) \mathbf{k}$, its limits, derivative,
and integrals are defined as the following.
\begin{itemize}
\item $\displaystyle \lim_{t \to a} \mathbf{r}(t) = 
    \left( \lim_{t \to a} f_1(t) \right) \mathbf{i} +
    \left( \lim_{t \to a} f_2(t) \right) \mathbf{j} +
    \left( \lim_{t \to a} f_3(t) \right) \mathbf{k}$
\item $\mathbf{r}'(t) = f_1'(t) \mathbf{i} + f_2'(t) \mathbf{j} +
    f_3'(t) \mathbf{k}$
\item $\displaystyle \int\mathbf{r}(t)\, dt =
    \int f_1(t)\, dt\ \mathbf{i} +
    \int f_2(t)\, dt\ \mathbf{j} +
    \int f_3(t)\, dt\ \mathbf{k}$
\item $\displaystyle \int_{a}^{b} \mathbf{r}(t)\, dt =
    \int_{a}^{b} f_1(t)\, dt\ \mathbf{i} +
    \int_{a}^{b} f_2(t)\, dt\ \mathbf{j} +
    \int_{a}^{b} f_3(t)\, dt\ \mathbf{k}$
\end{itemize}
\end{important}

Continuing with the definitions, it is natural to notice that the
limit, derivatives, and integrals exist if and only if they
exist for the component functions. In other words, $\mathbf{r}$
is continuous at $a$ if and only if its component functions are
continuous at $a$.

\begin{important}
Using the definition and properties of derivatives and integrals for
real-valued functions, it is self-evident that the following rules
also hold for vector-valued functions $\mathbf{r}$ under the
appropriate conditions.
\begin{itemize}
\item $(k\mathbf{r})' = k\mathbf{r}'$
\item $(\mathbf{r_1} \pm \mathbf{r_2})'=
    \mathbf{r_1}' \pm \mathbf{r_2}'$
\item For real-valued function $f$, $(f \cdot \mathbf{r})' =
    f' \cdot \mathbf{r} + f \cdot \mathbf{r}'$
\item $\displaystyle \int k\mathbf{r}(t)\, dt =
    k \int \mathbf{r}(t)\, dt$
\item $\displaystyle \int \mathbf{r_1}(t) \pm \mathbf{r_2}(t)\, dt =
    \int \mathbf{r_1}(t)\, dt \pm \int \mathbf{r_2}(t)\, dt$
\item For $\mathbf{R}'(t) = \mathbf{r}(t)$ satisfying for
$t \in [a, b]$, $\displaystyle \int_{a}^{b} \mathbf{r}(t)\, dt =
    \mathbf{R}(b) - \mathbf{R}(a)$
\end{itemize}
\end{important}

Before concluding this section of the note, let's discuss how we
could geometrically interpret derivatives of vector-valued functions.

First, consider the following equations.
\begin{align*}
    &\quad\ \mathbf{r}'(t) \\
    &= f_1'(t) \mathbf{i} + f_2'(t) \mathbf{j} + f_3'(t) \mathbf{k} \\
    &= \left(
        \lim_{h \to 0} \frac{f_1(t + h) - f_1(t)}{h}
    \right) \mathbf{i} + \left(
        \lim_{h \to 0} \frac{f_2(t + h) - f_2(t)}{h}
    \right) \mathbf{j} \\
    &\qquad\qquad\qquad\qquad\qquad\qquad\qquad\qquad\quad+ \left(
        \lim_{h \to 0} \frac{f_3(t + h) - f_3(t)}{h}
    \right) \mathbf{k} \\
    &= \lim_{h \to 0} \frac{f_1(t + h) \mathbf{i} +
        f_2(t + h) \mathbf{j} + f_3(t + h) \mathbf{k} - \left(
            f_1(t) \mathbf{i} + f_2(t) \mathbf{j} + f_3(t) \mathbf{k}
        \right)
    }{h} \\
    &= \lim_{h \to 0} \frac{\mathbf{r}(t + h) - \mathbf{r}(t)}{h}
\end{align*}

From our equation, it is evident that $\mathbf{r}'(t)$ is getting
more parallel with line tangent to the curve at the terminal point
of $\mathbf{r}(t)$ as $h \to 0$. In other words, $\mathbf{r}'(t)$
is a vector that is tangent to the curve at $r(t)$ for some value $t$.
Such vectors are called \textbf{tangent vector} and the line that
passes through tangent vectors is known as \textbf{tangent line}.
Let's take a look at an example.

\begin{problem}{}{}
Given $\mathbf{r}(t) = \langle \frac{t^2}{2}, \sin t, e^t \rangle$,
find the tangent vector and tangent line at $t = \pi$.
\end{problem}
\begin{solution}
First, notice that $\mathbf{r}'(t) = \langle t, \cos t, e^t \rangle$.
Therefore, $\mathbf{r}'(\pi) = \langle \pi, -1, e^\pi \rangle$,
which is the tangent vector.

Continuing, note that $\mathbf{r}(\pi) = \langle \frac{\pi^2}{2}, 0,
e^\pi \rangle$. Therefore, the tangent line can be represented as
$L(t) = \langle \frac{\pi^2}{2}, 0, e^\pi \rangle + t \langle \pi,
-1, e^\pi \rangle$.
\end{solution}

Before concluding this section of the note, consider the following
theorem.

\begin{theorem}{}{Derivates for Dot and Cross Products}
Given vector-valued functions $\mathbf{r_1}$ and $\mathbf{r_2}$,
the following equations hold.
\begin{itemize}
\item $(\mathbf{r_1} \cdot \mathbf{r_2})' = \mathbf{r_1}' \cdot
    \mathbf{r_2} + \mathbf{r_1} \cdot \mathbf{r_2}'$
\item $(\mathbf{r_1} \times \mathbf{r_2})' = \mathbf{r_1}' \times
    \mathbf{r_2} + \mathbf{r_1} \times \mathbf{r_2}'$
\end{itemize}
\end{theorem}
\begin{proof}
For the first equation, it suffices to show for vectors in 3-space
without loss of generality. Define $\mathbf{r_1} = \langle f_1, g_1,
h_1 \rangle$ and $\mathbf{r_2} = \langle f_2, g_2, h_2 \rangle$.
By definition, $\mathbf{r_1} \cdot \mathbf{r_2} = f_1f_2 + g_1g_2 +
h_1h_2$. Therefore, $(\mathbf{r_1} \cdot \mathbf{r_2})' =
(f_1'f_2 + g_1'g_2 + h_1'h_2) + (f_1f_2' + g_1g_2' + h_1h_2')$,
which is equal to $\langle f_1', g_1', \linebreak h_1' \rangle \cdot
\langle f_2, g_2, h_2 \rangle +  \langle f_1, g_1, h_1 \rangle \cdot
\langle f_2', g_2', h_2' \rangle$. Thus, the first equation holds.

Note that the second equation holds for $\mathbf{r_1}$ and
$\mathbf{r_2}$ in 3-space. Therefore, define $\mathbf{r_1} =
\langle f_1, g_1, h_1 \rangle$ and $\mathbf{r_2} =
\langle f_2, g_2, h_2 \rangle$. Consider the following equations.
\begin{align*}
    \mathbf{r_1} \times \mathbf{r_2}
    &= \langle
        g_1h_2 - h_1g_2, h_1f_2 - f_1h_2, f_1g_2 - g_1f_2
    \rangle \\
    (\mathbf{r_1} \times \mathbf{r_2})'
    &= \langle
        g_1'h_2 - h_1'g_2 + g_1h_2' - h_1g_2',
        h_1'f_2 - f_1'h_2 \\
    &\qquad\qquad\qquad\quad+ h_1f_2' - f_1h_2',
        f_1'g_2 - g_1'f_2 + f_1g_2' - g_1f_2' \rangle \\
    &= \langle
        g_1'h_2 - h_1'g_2, h_1'f_2 - f_1'h_2, f_1'g_2 - g_1'f_2
    \rangle \\
    \displaybreak
    &\qquad\qquad\qquad\quad+ \langle
        g_1h_2' - h_1g_2', h_1f_2' - f_1h_2', f_1g_2' - g_1f_2'
    \rangle \\
    &= \mathbf{r_1}' \times \mathbf{r_2} +
        \mathbf{r_1} \times \mathbf{r_2}'
\end{align*}
Thus, the theorem holds.
\end{proof}

One fun fact that we can notice from the first equation is that
if we assume $\mathbf{r}(t) \cdot \mathbf{r}(t) = a$ for some
constant $a$, then $\mathbf{r}'(t) \cdot \mathbf{r}(t) + \mathbf{r}(t)
\cdot \mathbf{r}'(t) = 0$ and $\mathbf{r}(t) \cdot \mathbf{r}'(t)
= 0$. In other words, $\mathbf{r}(t)$ and $\mathbf{r}'(t)$ are
orthogonal! This is it for this section, and let's discuss more on
parametrization in the next section.

\end{document}


